\section{Nef bases and extension of sections}
\label{V:sec-nef}

The nef-reduction arguments in Sections~\ref{V:sec-transfer} and
\ref{V:sec-invariant-field} start from twisted anticanonical sections.
Their common construction is a diagram
\[
 X\xleftarrow{p}W\xrightarrow{f}Y,
 \qquad ap^*L\sim D+f^*B,
\]
where $L=-K_X$, $a$ is a positive integer, $W$ is normal and
equidimensional over the smooth base $Y$, and $p$ is birational.
The divisor $D$ is effective Cartier, and above each prime of $Y$ it has
coefficient zero on at least one component.  The line bundle $B$ is
the part whose positivity must be established.

On a smooth resolution of $W$, a rank-one hypothesis for every relative
adjoint twist identifies the relevant sections with powers of the section
defining $D$, multiplied by the Jacobian section.  The bounded maximum
metric lemma below states this hypothesis in ordinary and invariant
forms.  It shows that the fiber maximum of the norm of the section defining $D$
gives a locally bounded semipositive metric on $B$, hence makes $B$
nef.  Nef reduction then identifies its numerical null directions, and
irregularity zero turns the resulting numerical descent into rational
linear descent.

The two applications finish differently.  For a prescribed arc,
Section~\ref{V:sec-transfer} proves $v(D)=0$ and extends first jets at
one point of the reduced base; this makes its descended line big and
yields arbitrarily small normalized arc orders.  Starting instead with
invariant twisted sections, Section~\ref{V:sec-invariant-field} extends
one nonzero invariant class in a quotient of multiplier ideals and
obtains an invariant anticanonical section.  We now prove the model,
metric, descent and curvature-concentration statements shared by these
constructions.

\subsection{Models and connected fibers}

The relative section must be normalized over every base divisor.
For that purpose we use a normal equidimensional model.  Its smooth
resolution is used for integration, and is not required to remain
equidimensional.

\begin{lemma}[Equidimensional models]\label{V:flat-model}
Let $X\dashrightarrow Y$ be a dominant rational map of integral projective
varieties over $\mathbb C$.  There are a smooth projective birational model
$Y'$ of $Y$ and an integral projective flat family $X_0\to Y'$ with a
birational morphism to the graph of the original map.  Every irreducible
component of every fiber has dimension $\dim X-\dim Y$.  The normalization
$V$ of $X_0$ is equidimensional over $Y'$; in particular, every vertical
prime divisor of $V$ maps onto a prime divisor of $Y'$.  If
$\mathbb C(Y)$ is algebraically closed in $\mathbb C(X)$, the morphism
from $V$, and from any smooth resolution of $V$, to $Y'$ has connected
fibers.  If a torus acts on $X$ and the rational map is invariant, these
constructions and a projective resolution of $V$ can be made equivariant,
with trivial action on the base.
\end{lemma}

\begin{proof}
Embed the projective graph over $Y$ in $\mathbb P^N\times Y$.
Over the flat open set the fibers define a map to the projective
Hilbert scheme \cite[Theorem~3.2 and the Hilbert-scheme specialization
after Proposition~3.8]{Grothendieck1961Hilbert}.
Resolve the closure of its graph, also resolving the base, and pull back the
universal family.  Flatness over an integral base embeds each affine
coordinate ring in its generic localization.  The generic fiber is
integral, so these coordinate rings have neither nilpotents nor zero
divisors.  Thus the pulled-back family is integral and is the closure of
the original family over the flat open set.  Its map to the graph is
birational.

Constancy of the Hilbert polynomial bounds the dimension of every fiber
above by the generic relative dimension.  At any fiber component the
height bound for the ideal generated by local parameters of the smooth
base gives the opposite bound.  This proves purity of fiber dimension.
Normalization is finite, so fiber dimensions and dimensions of images of
components are preserved.  A prime divisor mapping into codimension at
least two in $Y'$ would have dimension at most
$\dim Y'-2+(\dim X-\dim Y')=\dim X-2$, a contradiction.

For connectedness apply Stein factorization to the normal total space.
Its finite intermediate base has function field the algebraic closure
of $\mathbb C(Y')$ in $\mathbb C(X)$; the asserted field condition and
normality of $Y'$ make the finite birational morphism an isomorphism.
The same reasoning applies after resolution.  In the equivariant case,
the family over the original open set is stable under the action.
Its closure and normalization are therefore stable, and equivariant
resolution in characteristic zero supplies the last assertion.
\end{proof}

\subsection{A bounded metric from all relative moments}
\label{V:sec-maximum}

The two nef-reduction arguments below construct a section of a
relative line bundle and normalize its orders over base divisors.
The purpose of this normalization is analytic: it prevents the
maximum of the section's norm on nearby fibers from tending to zero.
Together with rank one for every relative adjoint twist, it produces
a nef line bundle on the base.  Ordinary and invariant rank one
are separate hypotheses in the statement.

\begin{lemma}[Bounded maximum metric]\label{V:bounded-maximum}
Let $X,Y$ be smooth connected projective varieties, let
$L=-K_X$ have a smooth semipositive metric $h$, and consider
\[
 W\xrightarrow{p}X,\qquad W\xrightarrow{f}Y,
 \qquad Z\xrightarrow{r}W,
\]
where $W$ is normal integral projective, $p$ is a birational morphism,
$f$ is surjective and
equidimensional with connected general fiber, and $r$ is a smooth
projective resolution.  Write $\pi=pr$, $g=fr$ and
$E=K_Z-\pi^*K_X\ge0$.  Let $B$ be a line bundle on $Y$, let $a\in\mathbb Z_{>0}$,
and suppose
\[
 0\ne\tau\in H^0(W,ap^*L-f^*B)
\]
has order zero on at least one component above every prime divisor
of $Y$.  Assume that, for every integer $l\ge1$, the ordinary
generic adjoint space for
$K_{Z/Y}+(al+1)\pi^*L$ has dimension one and is generated, in
base frames, by $s_E(r^*\tau)^l$.
Then $B$ has a semipositive metric with locally bounded weights.
On a smooth connected-fiber open set its weight in a frame $\beta$
is the continuous plurisubharmonic function
\begin{equation}\label{V:sup-weight}
 \varphi_\beta(y)=-\log\max_{Z_y}
       |r^*\tau\,g^*\beta|^2_{(\pi^*h)^a}.
\end{equation}
In particular, $B$ is nef.

The same conclusion holds with invariant generic rank one in place
of ordinary rank one, provided an algebraic torus acts compatibly
on the diagram with trivial action on $Y$, $L$ has its natural
linearization, $B$ has trivial linearization, $\tau$ is invariant,
and $h$ is invariant under the maximal compact torus.
\end{lemma}

\begin{proof}
The point-base case is immediate, so assume $\dim Y>0$.
The generic fiber is connected; the fibers over a sufficiently
small smooth-submersion open set $U$ are therefore connected and
smooth, hence irreducible.  Shrink $U$ so that neither $s_E$ nor
$r^*\tau$ vanishes identically on any fiber.

Choose frames $\beta$ of $B$ and $\kappa$ of $-K_Y$ on a base
chart.  For each $l\ge1$, the relative adjoint section
\[
 u_l=s_Eg^*\kappa\,(r^*\tau g^*\beta)^l
\]
generates the relevant generic space.  On its locally free
base-change locus, Berndtsson's smooth direct-image theorem
\cite[Theorem~1.2]{Berndtsson2009} applies to the smooth
semipositive coefficient $(al+1)\pi^*L$.  In the invariant case,
compact averaging is a holomorphic projection on the direct image.
It is orthogonal for the fiber integral metric, by change of
variables and invariance of the coefficient metric.  Its invariant
line is consequently also a metric quotient and inherits
semipositive curvature.  This is the only extra positivity step in
the invariant case.  If the relative dimension is zero, the
connected smooth fiber is one point and the same assertion follows
directly from the coefficient metric.

Put $b=|r^*\tau g^*\beta|^2_{(\pi^*h)^a}$, and let $d\mu_y$
be the smooth nonnegative density of $s_Eg^*\kappa$ with
coefficient $\pi^*L$.  The preceding positivity gives
plurisubharmonicity of
\[
 -\frac1l\log\int_{Z_y}b^l\,d\mu_y.
\]
The density has full support on each fiber in $U$.  Smooth proper
integration extends the displayed function across the possible
base-change exceptions for this particular $l$, preserving
plurisubharmonicity.  Thus no common base-change locus for infinitely
many $l$ is required.

Normalize each density to mass one.  Its $l$-th power mean of $b$
increases to the maximum.  The maximum is positive and continuous
under a proper smooth submersion, so this convergence is uniform
on compact subsets by monotonicity.  The logarithm of the mass,
divided by $l$, tends uniformly to zero there.  Passing to the
limit proves that \eqref{V:sup-weight} is plurisubharmonic on $U$.

It remains to extend the metric over $Y\setminus U$.  Regularity
of $\tau$ and properness bound $b$ above over compact base charts;
hence $\varphi_\beta$ is locally bounded below.  At a general point
of every boundary prime $Q$, choose a component $\Gamma$ above
$Q$ on which $\tau$ has order zero.  Normality makes $W$ smooth
at a general point of $\Gamma$.  We can choose this point so that
$\tau$ is nonzero and $\Gamma\to Q$ is submersive.  In local
coordinates a parameter for $Q$ pulls back to a unit times a
positive power of a parameter for $\Gamma$.  Taking a local root
of the unit, and completing with the tangential coordinates, shows
that $f$ is open there.  A neighborhood on which the norm of
$\tau$ is bounded below by a positive constant therefore covers
a base neighborhood.  Surjectivity of $Z\to W$ gives the same
lower bound for the fiber maximum.  Thus $\varphi_\beta$ is
locally bounded above near a general point of $Q$.

The removable-singularities theorem for plurisubharmonic functions
now extends the weights across the divisors, except over a closed
analytic set of codimension at least two.  The Hartogs extension
for plurisubharmonic functions removes that set as well.  More
explicitly, choose a small transverse disc whose boundary avoids
the missing set and parallel nearby discs with boundaries in a
fixed compact subset of the complement.  Outside a proper set of
parameters these discs miss the analytic set.  Their boundary
maximum gives a common upper bound, and the submean inequality
extends it to the other points of the domain.  The original lower
bounds persist under extension.  Finally,
\[
 \varphi_{u\beta}=\varphi_\beta-\log|u|^2
\]
for a nowhere-zero holomorphic function $u$, so the local weights
glue to a metric on $B$.  Restricting the bounded plurisubharmonic
weights to the normalization of any curve gives nonnegative
degree.  This proves nefness.
\end{proof}

The conclusion over the entire base is local boundedness.  Continuity
on the smooth locus is used for the power-mean limit; the extension
argument does not require continuity at every boundary point.

\subsection{Descending through the numerical null directions}

A nef divisor has degree zero on the general fibers of its nef
reduction.  The next argument upgrades this to every contracted
curve, which is the hypothesis needed for numerical pullback.
Irregularity zero is then used to turn that numerical identity into
a rational linear identity of line bundles.

\begin{lemma}[Descent of a nef divisor]\label{V:nef-descent}
Let $Y$ be smooth connected projective with $q(Y)=0$, and let $B$ be a
nef Cartier divisor on $Y$.  There are smooth connected projective
varieties $Y'$ and $Z$, a birational morphism $\rho:Y'\to Y$, a
surjective morphism $h:Y'\to Z$ with connected fibers, and a nef rational
Cartier divisor $L$ on $Z$ such that
\[
 \rho^*B\sim_{\mathbb Q}h^*L.
\]
If $\dim Z>0$, every irreducible curve through a very general point of
$Z$ has positive $L$-degree.  After multiplying $B$ and $L$ by one
positive integer, $L$ is Cartier and the displayed equivalence is linear.
\end{lemma}

\begin{proof}
Apply the nef reduction theorem of Bauer, Campana, Eckl, Kebekus,
Peternell, Rams, Szemberg and Wotzlaw
\cite[Theorem~2.1]{BCEKPRSW2002}.  It gives an almost holomorphic
dominant map $Y\dashrightarrow Z_0$ with connected general fibers,
numerical triviality of $B$ on its compact general fibers, and positive
$B$-degree for every noncontracted curve through a very general point
of $Y$.  Apply Lemma~\ref{V:flat-model} to its graph, over a smooth
projective birational model $Z$ of $Z_0$.  Write $T\to Z$ for the flat
family and $B_T$ for the pullback of $B$.

We claim that $B_T$ is numerically trivial on every curve in every fiber.
Almost holomorphicity identifies sufficiently general graph fibers with
compact fibers in the domain of the original map.  Indeed the complement
of that domain has smaller relative dimension over a general base point;
the compact fiber in the domain is already closed, and so acquires no
additional points in the graph.  Thus $B_T$ is numerically trivial on
a general fiber.

Let $d$ be the relative dimension.  If $d=0$ there are no contracted
curves on $T$.  Suppose $d>0$ and choose an ample Cartier divisor
$H$ on $T$.
Both $H$ and $H+B_T$ are relatively ample.  Their top degrees on the
fundamental fiber cycles are constant, by flatness and the leading
coefficients of their Hilbert polynomials, and agree on the general
fiber.  On any fiber the binomial expansion of their difference is a
sum of nonnegative mixed intersections.  Since the fiber has pure
dimension $d$, every component $R$ occurs with positive multiplicity;
we obtain
\[
 B_T\cdot H^{d-1}\cdot R=0
 \qquad\hbox{for every fiber component }R.
\]
If $d=1$, every integral curve $C\subset R$ equals $R$, so its degree
is already zero. For $d\ge2$, choose $d-1$ sufficiently high ample
hypersurfaces through a prescribed curve $C\subset R$. Global generation
of its twisted ideal permits successive general choices avoiding all
existing components not contained in $C$, so the final intersection is
proper and one-dimensional. At the generic point of $C$, their local
equations form a system of parameters; their intersection multiplicity
there is positive. Thus the resulting effective curve cycle contains
$C$ with positive coefficient, also when $R$ is singular.
Its $B_T$-degree is
zero, and nefness makes every component degree nonnegative.  Hence
$B_T\cdot C=0$, as claimed.

Normalize $T$ and resolve it to $Y'$, with maps $\rho:Y'\to Y$ and
$h:Y'\to Z$.  A curve contracted by $h$ either maps to a curve in a
fiber of $T$ or maps to a point, so $\rho^*B$ has degree zero on every
such curve.  The fibers of $h$ are connected, by Stein factorization
and the connected general fibers of the nef reduction.  Lehmann's
numerical pullback theorem \cite[Theorem~1.2]{Lehmann2015} applies:
the divisor $\rho^*B$ is pseudoeffective, is numerically trivial on
all contracted curves, and the smooth base $Z$ is $\mathbb Q$-factorial.
It gives $\rho^*B\equiv h^*L$ for a real numerical divisor class $L$.
The pullback map on numerical divisor spaces is defined over
$\mathbb Q$, so the linear equation with rational right-hand side
has a rational solution.  Choose such an $L$.  It is nef: for each
curve in $Z$, a projective slice of its inverse image gives a curve
dominating it, and the projection formula gives nonnegative degree.

Birational invariance gives $q(Y')=0$.  Numerically trivial integral
line bundles on a smooth projective variety belong to $\operatorname{Pic}^0$
up to torsion; see \cite[Definition~6.1 and Theorem~6.3]{Kleiman2005}.  Therefore the rational numerical equivalence just
obtained is a rational linear equivalence.  Clearing its torsion and
denominators gives the asserted integral scaling.

Finally, choose a very general $z\in Z$ and a point $y'\in h^{-1}(z)$
where $h$ is smooth and $\rho$ is an isomorphism, such that $\rho(y')$
lies outside the countable exceptional union in the nef reduction
theorem.  Such a point exists: for each proper closed exception, the
image of its complement contains a dense open subset of $Z$; on a very
general irreducible fiber, the resulting countably many proper closed
subsets cannot cover the fiber.  Given any curve through $z$, normalize
it and slice a component of its inverse image through $y'$ to obtain
a curve through $y'$ dominating it.  Smoothness at $y'$ permits a
nonvertical curve germ, so the slicing may be chosen to retain that
property.  Its image in $Y$ is not contracted by the original nef
reduction.  The positivity part of that theorem, followed by the
projection formula, proves that the given curve has positive
$L$-degree.
\end{proof}

The reduced base now has positive degree on every curve through a
very general point.  We next turn this qualitative condition into
an arbitrarily large Seshadri bound after adding a fixed small ample
term.  This will give an isolated logarithmic pole strong enough to
detect either a first jet or a nonzero quotient class.

We use the following metric conventions in this extension argument.
If a singular Hermitian metric has squared frame norm $h=e^{-\varphi}$,
its multiplier ideal $\mathcal I(h)$, also denoted $\mathcal J(h)$,
consists of holomorphic germs $u$ for which $|u|^2e^{-\varphi}$ is
locally integrable. The metric $he^{-\psi}$ has weight
$\varphi+\psi$ and curvature $i\Theta_h+i\partial\bar\partial\psi$.
A quasi-plurisubharmonic function has \emph{neat analytic singularities}
if locally it has the form $c\log\sum_j|f_j|^2+u$, where $c>0$,
the $f_j$ are holomorphic and $u$ is smooth.

\begin{lemma}[Concentrating curvature at one base point]
\label{V:sesha}
Let $Z$ be smooth projective of dimension $b>0$, let $L$ be a nef
Cartier divisor, and suppose every curve through a very general point
of $Z$ has positive $L$-degree.  Fix a very ample Cartier divisor $H$
and rational $\epsilon>0$.  For every $c>0$, all sufficiently large
integers $k$ make
\[
 J_k=\tfrac12(kL+\epsilon H)
\]
an ample rational divisor with Seshadri constant greater than $c$ at
a very general point.  At any such point $z$ one can choose a smooth
positive metric of curvature $\eta$ on $J_k$ and a quasi-plurisubharmonic
function $\psi$ with neat analytic singularities, smooth away from $z$,
such that
\[
 \eta+i\partial\bar\partial\psi\ge0,
 \qquad \psi(w)\le c\log|w|^2+O(1)
\]
in coordinates centered at $z$.
\end{lemma}

\begin{proof}
Fix the countable union of proper closed subsets outside which the
curve positivity holds.  If an integral $d$-dimensional subvariety
$U$ is not contained in this union, then
$L\cdot H^{d-1}\cdot U$ is a positive integer.  Otherwise nefness
and the cutting-through-a-curve argument in the proof of
Lemma~\ref{V:nef-descent} would make $L$ numerically trivial on
$U$, contrary to a curve in $U$ through a point outside the union.
The mixed term in the binomial expansion consequently gives
\[
 J_k^d\cdot U\ge 2^{-d}d k\epsilon^{d-1}.
\]
Choose $\alpha>c$ and then $k$ large enough that this is at least
$(d\alpha)^d$ for every $1\le d\le b$.  The theorem of
Ein--K\"uchle--Lazarsfeld \cite[Theorem~3.1]{EKL1995}, applied after
clearing rational denominators, gives the required Seshadri bound.

Choose a rational number $c_0$ strictly between $c$ and the Seshadri
constant.  On the blowup of $z$, the pullback of $J_k$ minus $c_0$
times the exceptional divisor is ample.  To see ampleness rather than
mere nefness, choose $c_0<c'<\epsilon(J_k;z)$; a sufficiently small positive
exceptional subtraction is ample, and the desired divisor is a
positive convex combination of that ample divisor and the nef divisor
with coefficient $c'$.  For a sufficiently divisible large integer
$p$, a basis of sections of $pJ_k$ vanishing to order at least $pc_0$
at $z$ therefore has no common zero away from $z$.  If these sections
are $s_j$, put
\[
 \psi=\frac1p\log\sum_j|s_j|^2_{h_{J_k}^{p}}.
\]
The curvature inequality follows directly, and the pole with coefficient
$c_0>c$ gives the asserted upper bound after shrinking the coordinate
neighborhood.  For $\dim Z=1$ use $J_k-c_0z$ on the curve.
\end{proof}

